% ===== PRÉAMBULE CANONIQUE — à coller VERBATIM en tête de CHAQUE .tex =====
\documentclass[11pt,a4paper]{article}
\usepackage[utf8]{inputenc}\usepackage[T1]{fontenc}\usepackage{lmodern}
\usepackage[french]{babel}
\usepackage{amsmath,amssymb,mathtools}\usepackage{icomma}
\usepackage[version=4]{mhchem}          % \ce{...} : équations & formules
\usepackage{siunitx}\sisetup{locale=FR} % \qty{}{}, \si{}, \ang{}
\usepackage{chemfig}                    % \chemfig{...} : structures développées
\usepackage{xcolor}\usepackage{colortbl}
\usepackage{tikz}\usepackage{pgfplots}\pgfplotsset{compat=1.18}
\usepackage[most]{tcolorbox}\usepackage{enumitem}\usepackage{array,booktabs}
\usepackage[a4paper,margin=2.2cm]{geometry}
\usetikzlibrary{arrows.meta,calc,angles,quotes,positioning,patterns,backgrounds,shapes.geometric}
\definecolor{primary}{RGB}{222,49,99}\definecolor{primarydark}{RGB}{150,22,55}
\definecolor{defcol}{RGB}{0,114,178}\definecolor{methcol}{RGB}{52,92,148}
\definecolor{excol}{RGB}{0,135,90}\definecolor{attncol}{RGB}{200,40,55}
\tcbset{boxbase/.style={enhanced,breakable,boxrule=0.9pt,arc=2.5pt,
  left=3mm,right=3mm,top=2.3mm,bottom=2.3mm,fonttitle=\bfseries,coltitle=white}}
% --- boîtes TUTORIEL (noms EXACTS, ne pas renommer) ---
\newtcolorbox{cle}[1][]{boxbase,colback=defcol!6,colframe=defcol,
  title={\textsc{À retenir}\ifx\relax#1\relax\relax\else\textmd{ : \itshape #1}\fi}}
\newtcolorbox{methode}[1][]{boxbase,colback=methcol!7,colframe=methcol,sharp corners,
  title={\textsc{Méthode}\ifx\relax#1\relax\relax\else\textmd{ --- \itshape #1}\fi}}
\newtcolorbox{exemplebox}[1][]{boxbase,colback=excol!5,colframe=excol,
  title={\textsc{Exemple}\ifx\relax#1\relax\relax\else\textmd{ : \itshape #1}\fi}}
\newtcolorbox{piege}[1][]{boxbase,colback=attncol!6,colframe=attncol,
  title={\textsc{Piège}\ifx\relax#1\relax\relax\else\textmd{ : \itshape #1}\fi}}
\newtcolorbox{remarque}[1][]{boxbase,colback=black!4,colframe=black!45,
  title={\textsc{Remarque}\ifx\relax#1\relax\relax\else\textmd{ : \itshape #1}\fi}}
% --- boîtes EXERCICE / CORRIGÉ (noms EXACTS, auto-comptées) ---
\newtcolorbox[auto counter]{exo}[1][]{boxbase,colback=primary!4,colframe=primary,
  title={\textsc{Exercice}~\thetcbcounter\ifx\relax#1\relax\relax\else\textmd{ --- \itshape #1}\fi}}
\newtcolorbox[auto counter]{corrige}[1][]{boxbase,colback=excol!4,colframe=excol,
  title={\textsc{Corrigé}~\thetcbcounter\ifx\relax#1\relax\relax\else\textmd{ --- \itshape #1}\fi}}
\newcommand{\R}{\mathbb{R}}\newcommand{\N}{\mathbb{N}}\newcommand{\Z}{\mathbb{Z}}\newcommand{\ds}{\displaystyle}
% (un \usepackage{fancyhdr}+en-tête est possible ; il est ignoré côté web)
% =========================================================================
\begin{document}
\sisetup{per-mode=symbol}
\begin{exo}[une pesée dans toutes les unités]
Une microbalance affiche $m = 4,50\times10^{-3}\ \si{\gram}$. Exprime cette \emph{même} masse en
milligrammes, microgrammes et kilogrammes, en notation scientifique et \emph{sans} changer le nombre de
chiffres significatifs.
\begin{enumerate}[label=\alph*)]
  \item en \si{\milli\gram}
  \item en \si{\micro\gram}
  \item en \si{\kilo\gram}
\end{enumerate}
\end{exo}

\begin{exo}[un même calcul, trois précisions]
Le calcul brut du volume d'un cube donne $V = 15,625\ \si{\centi\metre\cubed}$. Arrondis ce résultat :
\begin{enumerate}[label=\alph*)]
  \item à $4$ chiffres significatifs
  \item à $3$ chiffres significatifs
  \item à $2$ chiffres significatifs
\end{enumerate}
\end{exo}

\begin{exo}[énergie de combustion par kilogramme]
La combustion de $\qty{1}{\gram}$ d'essence libère $4,307\times10^{4}\ \si{\joule}$. On rappelle que
$\qty{1}{\gram}=10^{-3}\ \si{\kilo\gram}$ (facteur exact). Exprime l'énergie libérée \emph{par
kilogramme}, en notation scientifique et à $3$ chiffres significatifs.
\end{exo}

\begin{exo}[vrai ou faux, justifié]
Indique si chaque affirmation est \textbf{vraie} ou \textbf{fausse}, et corrige-la si besoin.
\begin{enumerate}[label=\alph*)]
  \item «~Arrondir $0,7999$ à $2$ CS donne $0,80$.~»
  \item «~$2,5\times10^{-3}$ et $0,0025$ sont deux écritures du même nombre.~»
  \item «~En notation scientifique, la mantisse peut valoir $12,5$.~»
  \item «~On peut arrondir à n'importe quelle étape d'un calcul sans changer le résultat final.~»
\end{enumerate}
\end{exo}

\end{document}
